\documentclass[a4paper,10pt]{article}
\usepackage{cv}

\hypersetup{pdftitle={陆知雨 - 简历}}

% 中文版：覆盖合作者命令
\renewcommand*{\with}[1]{\,（合作者：#1）}

\begin{document}

\name{陆知雨}

\begin{center}
  \small 中央财经大学财政学博士候选人
\end{center}
\vspace{-0.5em}

\contactblock{%
  中央财经大学\\
  中国财政发展协同创新中心%
}{%
  \begin{tabular}{@{}ll@{}}
    邮箱：    & \href{mailto:luzhiyu@email.cufe.edu.cn}{\texttt{luzhiyu@email.cufe.edu.cn}} \\[2pt]
    个人网站：& \href{https://zhiyulu.org}{\texttt{zhiyulu.org}} \\
  \end{tabular}%
}

% -------------------------------------------------------
\block{教育经历}
% -------------------------------------------------------

\begin{itemize}
\item 中央财经大学财政税务学院\\
  经济学学士，专业：财政学（2022年9月--2026年6月）\\
  核心课程：Python编程（99）、运筹学（100）、宏观经济学（97）、计量经济学（95）、数据科学导论（95）、博弈论（96）、组织经济学（92）、会计学（93）

\item 中央财经大学中国财政发展协同创新中心\\
  财政学博士研究生（2026年9月至今）

\end{itemize}

% -------------------------------------------------------
\block{研究兴趣}
% -------------------------------------------------------

\begin{itemize}
\item \textsb{研究领域：}政治经济学 \sep 公共经济学 \sep 创新经济学
\item \textsb{研究方法：}文本数据 \sep 机器学习 \sep 行政数据
\end{itemize}

% -------------------------------------------------------
\block{工作论文}
% -------------------------------------------------------

\begin{itemize}

\item Licensed to Diverge: Central Intervention and Local Policy Rewriting in China\with{姚东旻、徐傲晗}\\
  学术报告：南开大学

\item Hidden Injustice: Gender Inequality in Divorce Judgments and Judicial Delocalization\with{姚东旻、熊国豪、高洋、刘芳妍}\\
  会议报告：\textit{香樟经济学年会 2026}，上海财经大学，2026.05

\item Public Demand Expansion and Private Enterprise Dynamism: Evidence from Chinese Government Procurement\with{姚东旻、王子豪、张世宇}

\item 数字治理能否提升城市创业活力？——基于政府采购数据的证据\with{阮睿、沈昕毅、孙宇辰}\\
  会议报告：\textit{中国组织经济学研讨会}，辽宁大学，2025.05；\textit{香樟数字经济论坛}，河南财经政法大学，2025

\item Cultural Differences in Large Language Models: An Ultimatum Game Study across Models, Languages, and Identity Conditions\with{姚东旻、董鸿儒}

\end{itemize}

% -------------------------------------------------------
\block{科研经历}
% -------------------------------------------------------

\begin{itemize}

\item 财政政策与货币政策有效协调配合的体制机制研究 \sep 助研 \sep 2025至今\\
  {\color{gray}\textit{国家社科基金重大项目，中央财经大学}}

\item 基于大数据的经济政策成本与收益评价通用模型研究 \sep 助研 \sep 2025至今\\
  {\color{gray}\textit{国务院发展研究中心课题，中央财经大学}}

\item 主要经济体财税改革和财税政策动态研究 \sep 助研 \sep 2025至今\\
  {\color{gray}\textit{财政部政策研究室课题，中央财经大学}}

\item 分领域中央与地方财政事权和支出责任划分改革实效评估 \sep 助研 \sep 2024至今\\
  {\color{gray}\textit{财政部政策研究室课题，中央财经大学}}

\item 政府预决算视角下提升我国国家创新体系整体效能的财政体制与政策研究 \sep 助研 \sep 2020--2025\\
  {\color{gray}\textit{国家社科基金重点项目，中央财经大学}}

\item 中国土地财政与城市财政韧性研究 \sep 助研 \sep 2024.01--2025.04\\
  {\color{gray}\textit{北京大学—林肯研究院城市发展与土地政策研究中心}}

\end{itemize}

% -------------------------------------------------------
\block{软件与工具}
% -------------------------------------------------------

\begin{itemize}

\item \href{https://pypi.org/project/pybdiff/}{pybdiff}\\
  Stata \texttt{bdiff} 命令的 Python 复现。通过置换检验、Bootstrap 或 Wald 检验，检验两组回归系数是否存在显著差异，支持并行计算。可通过 \texttt{pip install pybdiff} 安装。

\item \href{https://pypi.org/project/pyshapley2/}{pyshapley2}\\
  Stata \texttt{shapley2} 命令的 Python 复现。计算任意回归拟合统计量的 Shapley--Owen 分解，支持自定义变量分组与并行计算。可通过 \texttt{pip install pyshapley2} 安装。

\item \href{https://pypi.org/project/pybacondecomp/}{pybacondecomp}\\
  Stata \texttt{bacondecomp} 命令的 Python 复现。实现 Goodman--Bacon（2021）对双向固定效应 DiD 估计量的 $2\times2$ 分解，支持并行计算。可通过 \texttt{pip install pybacondecomp} 安装。

\item \href{https://jrnldaily.com}{jrnldaily.com}\\
  经济学文献自动聚合平台，每日收录最新工作论文与期刊文章，提供 AI 摘要生成与跨语言检索功能。

\end{itemize}

% -------------------------------------------------------
\block{获奖情况}
% -------------------------------------------------------

\begin{itemize}

\item 全面发展奖学金 \sep 中央财经大学 \sep 2025
\item 学术科研与创新优秀奖学金 \sep 中央财经大学 \sep 2025，2026
\item 学术德育优秀奖学金 \sep 中央财经大学 \sep 2025
\item 全国大学生公共经济与政策论文大赛三等奖 \sep 2024
\item 大学生创新创业训练计划：国家级优秀结项 2026 \sep 北京市优秀结项 2024 \sep 北京市级结项 2025

\end{itemize}

\bigskip
\noindent{\footnotesize\textcolor{gray}{版本：2026年11月}}

\end{document}
